\documentclass{article}

% NeurIPS 2026 workshop style, double-blind submission.
% The [dblblindworkshop] option gives the anonymized workshop format; do NOT add
% [final] (that de-anonymizes for camera-ready). neurips_2026.sty is bundled here.
\usepackage[dblblindworkshop]{neurips_2026}

\usepackage[utf8]{inputenc}
\usepackage[T1]{fontenc}
\usepackage{hyperref}
\usepackage{url}
\usepackage{booktabs}
\usepackage{amsmath}
\usepackage{graphicx}

\title{Silent Grounding Failures: A Formal-Complexity Account of\\
Parser Reliability for Agentic Molecular Science}

% Required by the NeurIPS 2026 workshop template.
\workshoptitle{ML4Molecules: Agentic Systems for Molecular Sciences (NeurIPS 2026)}

\author{%
  Anonymous Author(s)\\
  Affiliation\\
  Address\\
  \texttt{email}\\
}

\begin{document}

\maketitle

\begin{abstract}
Agentic ``co-scientists'' in the molecular sciences reason over the \emph{results}
of simulations---geometries, energies, spectra, crystal structures---which they
must first extract from the raw output files of quantum-chemistry and
materials codes. This grounding step is universally assumed to work. We test the
assumption. Exposing four widely used community parsers as agent tools, we
benchmark them on 1{,}104 real files from 52 entries of a public materials
database spanning 16 simulation codes, scoring by whether an atomic structure was
actually recovered rather than by whether an exception was avoided. Reliability
is poor and, worse, often \emph{silent}: on the single file each entry designates
as its primary output, no parser exceeds 24\% success, an oracle picking the best
parser per file reaches only 45\%, and one parser returns empty-but-valid objects
without error on 98\% of inputs---indistinguishable from success to a calling
agent. We then show these failures are not arbitrary. Classifying each format by
the grammatical complexity its extraction demands (a coarse Chomsky-hierarchy
ordinal), success falls monotonically with complexity: structured formats are
grounded 100\% of the time, free-form logs 43\%, cross-referential multi-file
outputs 0\%. We frame this as a \emph{power gap}: a fixed parser fails when the
format's complexity exceeds its parsing power, which also predicts when an
agent should stop calling tools and synthesize its own parser. Our results are a
cautionary, rigorously-measured baseline for the data-grounding layer beneath
agentic molecular science.
\end{abstract}

\section{Introduction}

Recent work casts large language model agents as co-scientists that plan
syntheses, orchestrate domain tools, and close experimental loops. Careful
benchmarking keeps exposing fragility at surprisingly basic tasks---tool-augmented
agents near 50\% on chemical cost estimation, single-cell foundation models that
do not beat linear baselines. We add a data point beneath all of these: the point
where an agent turns a simulation's \emph{output file} into the structured numbers
it will reason over.

Every agentic molecular workflow that consumes the results of a computation---to
compare reaction energies, screen a crystal, fit an interatomic potential, or
verify a geometry---must first parse those results out of a code's output. The
community has mature parsers for this, and agent frameworks increasingly wrap them
as callable tools. The implicit assumption is that these tools ``just work.'' When
they do not, the failure is a scientific one: the agent reasons over an empty or
wrong structure and reaches a confidently wrong molecular conclusion, with nothing
in the pipeline to signal that the grounding step failed.

We make two contributions. First (\S\ref{sec:bench}--\ref{sec:results}), a
rigorous reliability benchmark of the grounding layer on real, heterogeneous
simulation output, using a field-level success criterion (\emph{did we recover an
atomic structure?}) rather than the weaker ``no exception raised.'' The picture is
poor and frequently silent. Second (\S\ref{sec:complexity}), we show the failures
are \emph{predictable}: classifying each source format by the grammatical
complexity its extraction demands turns ``parsers fail unpredictably'' into a
monotone law, and recasts the coding agent from a caller of tools into the
mechanism that must climb the complexity hierarchy when the fixed tools are
under-powered.

\section{The grounding layer and how we benchmark it}
\label{sec:bench}

\paragraph{Corpus.} We draw 52 entries from a public computational-materials
database, spanning 16 simulation codes (VASP, Gaussian, ORCA, GAMESS, exciting,
CASTEP, CP2K, ABINIT, Crystal, FHI-aims, GPAW, ONETEP, NWChem, LAMMPS, Octopus,
WIEN2k), comprising 1{,}104 raw files. Each entry carries an authoritative
\emph{mainfile}: the single primary-output file the database designates as the
entry's result. The corpus, program metadata, and mainfiles are frozen for
reproducibility.

\paragraph{Parsers as tools.} We evaluate four parsers exposed through a uniform
\texttt{parse(path)$\rightarrow$model} tool interface: two general community
libraries for atomistic and quantum-chemistry output, a general
input/output library with format-specific grammars, and a custom regular-expression
parser for one code. Each targets a common set of molecular quantities: formula,
charge, multiplicity, atom count, Cartesian geometry, total/SCF energy,
vibrational frequencies and IR intensities, and atomic charges.

\paragraph{Two denominators.} We report every parser against (i) \emph{all files}
in an entry---the naive ``point the tool at everything'' regime---and (ii)
\emph{mainfile only}, the primary output. The gap separates ``failed on a file it
never claimed to handle'' from genuine failure on the intended input.

\paragraph{Field-level success, not ``no exception.''} Every parser returns an
object whose fields are optional, so a parser can return an essentially empty
object without raising. We therefore define \textbf{success} as recovering an
atomic structure (atom count and coordinates), and separately track energy
recovery. Each attempt is labelled \texttt{success}, \texttt{empty} (no exception,
no structure---a \emph{silent} failure), \texttt{error} (exception; type
recorded), or \texttt{timeout}. Parsers are never modified: their failures are the
measurement.

\section{The grounding layer is unreliable---and fails silently}
\label{sec:results}

\begin{table}[t]
\centering
\small
\begin{tabular}{lcc}
\toprule
Parser & All files & Mainfile only \\
\midrule
Atomistic (A)      & 3.9\% (43/1104)  & 23.5\% (12/51) \\
Quantum-chem (B)   & 15.4\% (170/1104)& 21.6\% (11/51) \\
Custom regex (C)   & 2.4\% (26/1104)  & 7.8\% (4/51)   \\
I/O library (D)    & 10.5\% (116/1104)& 5.9\% (3/51)   \\
\midrule
\textbf{Oracle (any parser)} & \textbf{20.5\% (226/1104)} & \textbf{45.1\% (23/51)} \\
\bottomrule
\end{tabular}
\caption{Structure-recovery success. Even on the designated primary output, no
single parser exceeds 24\%, and an oracle that picks the best parser per file
solves under half; on the remaining 55\% of mainfiles no parser recovers a
structure at all.}
\label{tab:reliability}
\end{table}

\begin{table}[t]
\centering
\small
\begin{tabular}{lcccc}
\toprule
Parser & Success & Empty (silent) & Error & Timeout \\
\midrule
Atomistic (A)    & 43  & 0    & 1061 & 0  \\
Quantum-chem (B) & 170 & 39   & 849  & 46 \\
Custom regex (C) & 26  & 1077 & 0    & 1  \\
I/O library (D)  & 116 & 163  & 825  & 0  \\
\bottomrule
\end{tabular}
\caption{Outcome breakdown (all files). The \emph{Empty} column is the hazard: the
custom parser raises no exceptions yet returns a structureless object on 1077/1104
inputs. A calling agent cannot distinguish these from success without inspecting
the payload.}
\label{tab:outcomes}
\end{table}

Table~\ref{tab:reliability} reports the headline. On the mainfile---the file the
database itself flags as the result, where ``it just works'' should hold most
strongly---no single parser exceeds 23.5\%, and the oracle reaches 45.1\%.

\paragraph{Failures are often silent.} Table~\ref{tab:outcomes} shows the more
dangerous pattern. Exceptions are catchable; a structureless object returned
\emph{without} an exception is not. The custom parser never raises and yet grounds
nothing on 98\% of files; the I/O library and quantum-chemistry parser do the same
on many others. An agent that guards only against exceptions---the common
case---will silently ingest nothing and reason on it.

\paragraph{Failure taxonomy.} The dominant explicit failure is \emph{format not
recognized}: for the atomistic parser, unknown-file-type (39\%), premature
end-of-iteration on truncated or malformed files (21\%), value errors (16\%), and
Unicode-decode errors on binary interleaving (13\%); for the quantum-chemistry
parser, 95\% of exceptions are its ``format unrecognized'' signal; for the I/O
library, 97\%. Recovery of a total energy is even rarer than structure---the
atomistic parser recovers an energy on 0 of 51 mainfiles, reading geometries but
not the energetics an agent usually wants.

\section{A formal-complexity account of the failures}
\label{sec:complexity}

These parsers are hand-written scanners over Fortran/C \texttt{write}-statement
output, not generated from grammars. Whether one succeeds should therefore depend
on the grammatical complexity extraction demands versus the class of languages the
parser can recognize. We make this precise with a coarse ordinal inspired by the
Chomsky hierarchy---a documented rubric, not a formal per-file classification:
\textbf{L0} grammar-defined/structured (XML, JSON, checkpoint/orbital files;
context-free); \textbf{L1} line-regular records (regular); \textbf{L2} free-form
logs with count-agreement (a header declares $N$ atoms, then $N$ lines follow;
repeated blocks where the last must be taken; mildly context-sensitive);
\textbf{L3} cross-referential or multi-file outputs (values defined in one
file/section and combined in another; context-sensitive).

\begin{table}[t]
\centering
\small
\begin{tabular}{lccccc}
\toprule
Level & A & B & C & D & Oracle \\
\midrule
L0 (structured)        & 100\% (4/4)  & 0\%          & 0\%         & 0\%          & \textbf{100\% (4/4)} \\
L2 (free-form log)     & 18.2\% (8/44)& 25.0\% (11/44)& 9.1\% (4/44)& 6.8\% (3/44) & \textbf{43.2\% (19/44)} \\
L3 (cross-referential) & 0\% (0/3)    & 0\%          & 0\%         & 0\%          & \textbf{0\% (0/3)} \\
\bottomrule
\end{tabular}
\caption{Mainfile structure-recovery re-cut by format complexity. Oracle success
falls monotonically as complexity rises: structured formats are grounded every
time, free-form logs under half, cross-referential outputs never.}
\label{tab:complexity}
\end{table}

\paragraph{The power gap.} Table~\ref{tab:complexity} re-cuts the mainfile results
by level. Oracle success is monotone: L0 100\%, L2 43\%, L3 0\%. This supports a
\emph{power-gap hypothesis}: a fixed parser fails when the source's level exceeds
its parsing-power ceiling (regex parsers are regular; stateful scanners handle
counting; only grammar-backed readers reach context-free structure). It also
reframes the coding agent: writing and running code is Turing-complete, so the
agent's role is precisely to \emph{climb} the hierarchy---to synthesize a stateful
or recursive parser when every fixed tool is under-powered. The same lens predicts
when the agent should do so (L3 or unsupported formats) versus when it wastes
effort re-implementing a tool that already covers the format (L0--L1).

\paragraph{Complexity is necessary, not sufficient---three axes.} We are careful
not to over-fit the hierarchy. Two large parts of the failure are orthogonal to
grammatical complexity. \emph{Identification:} the dominant explicit failure was
``format not recognized''---the parser never dispatched---which is a routing
problem, not a power problem; indeed L0 files in program-specific schemas that no
tool implements score 0\%, so context-free structure is necessary but not
sufficient. \emph{Noise:} binary interleaving, truncation, and version drift break
parsers on formats that are otherwise trivial. Extraction difficulty thus
decomposes into three axes---\emph{identify}, \emph{parse}, \emph{denoise}---of
which the Chomsky lens cleanly governs only the second. We also note a confound:
at $N{=}4$ (L0) and $N{=}3$ (L3), level is partly collinear with program in this
corpus, so the gradient is suggestive rather than causal.

\paragraph{Relation to prior work.} Grouping tasks by the Chomsky hierarchy to
predict neural generalization was established by \citet{deletang2023chomsky}, and a
line of expressivity results shows chain-of-thought and tool use lift transformers
toward Turing-completeness \citep{merrill2024cot,strobl2024survey}; we borrow this
to justify the agent-as-escalator view rather than to contribute to it. Extracting
structure from semi-structured machine text is the subject of automated log parsing
\citep{zhu2019logparsing}, but that line mines recurring templates for anomaly
detection, not conversion to a scientific target schema, and does not use a
complexity lens. LLM schema-extraction in science targets natural-language prose
\citep{dagdelen2024extraction}, a different regime from machine-emitted output
files. The parsers we evaluate \citep{cclib} and the hierarchical target schema
and per-code parser infrastructure of open materials databases
\citep{ghiringhelli2023metadata} are the components we stress-test. To our
knowledge, framing tool-layer reliability for output-file$\rightarrow$schema
conversion as a parsing-power/format-complexity gap is new.

\section{Implications for agentic molecular science}

Three takeaways for systems that ground on simulation output. (1) \emph{Guard on
content, not exceptions.} The silent-\texttt{empty} outcome shows exception
handling is insufficient; a grounding tool must assert that scientific content was
actually recovered before an agent reasons over it. (2) \emph{Route by complexity,
and expect a ceiling.} Success is strongly format-dependent and the winning parser
differs by domain, so a capability- and complexity-aware router over multiple
parsers is the right architecture---but Table~\ref{tab:reliability} bounds what
routing buys (45\% on primary files); the rest needs new parsers, format hints, or
agent-synthesized parsing. (3) \emph{Report the denominator.} Naive success rates
understate parsers dramatically (3.9\% vs.\ 23.5\% for parser A); evaluations of
tool reliability must state what set they score against.

\paragraph{Outlook.} The power-gap hypothesis makes a testable prediction about
\emph{agents}: giving a coding agent the source's complexity level and each tool's
power ceiling should reduce wasted wrong-parser calls and trigger correct
parser-from-scratch synthesis exactly where the fixed tools are under-powered. A
controlled ablation of such scaffolding is our immediate next step.

\section*{Limitations}

The corpus is a single snapshot of one materials database (52 entries, 16 codes);
success is defined by field \emph{presence}, not correctness against ground truth,
so a parser that recovers a \emph{wrong} structure counts as success; parsers are
run with defaults and no format hints, which is an agent's default but not an
expert's; the complexity levels are a coarse, documented ordinal rather than a
formal classification, and are partly confounded with program at small $N$; and we
benchmark the tools, not a live agent loop---the agentic ablation is future work.

{\small
\bibliographystyle{plainnat}
\bibliography{refs}
}

\end{document}
